(1) 基尔霍夫定律为 \(L\dot I+RI=s(t)V\)，其中前半周期 \(s=+1\)、后半周期 \(s=-1\)。每段的通解均为 \(I(t)=I_{\rm eq}+[I(t_0)-I_{\rm eq}]e^{-(t-t_0)/\tau}\)，\(\tau=L/R\)，其中 \(I_{\rm eq}=\pm V/R\)。瞬间切换电压时电感电流连续。
(2) 令 \(h=e^{-P/(2\tau)}\)，在第一段末 \(I_1=V/R+(I_0-V/R)h\)；第二段末返回 \(I_0=-V/R+(I_1+V/R)h\)。解得 \(I_0=-(V/R)(1-h)/(1+h)\)、\(I_1=(V/R)(1-h)/(1+h)\)，具有反半周期对称性。
(3) 在 \(0\le t<P/2\)，\(I(t)=V/R+[I_0-V/R]e^{-t/\tau}\)；后半段 \(I(t)=-V/R+[I_1+V/R]e^{-(t-P/2)/\tau}\)。由对称关系 \(I(t+P/2)=-I(t)\) 可直接得 \(\langle I\rangle_P=0\)，但电流平方平均值非零，仍存在正焦耳热。
(4) 电源瞬时功率为 \(VI s(t)\)，乘回路方程以 \(I\) 得 \(VIs=RI^2+d(LI^2/2)/dt\)。稳态一周期首末电流相同，电感储能变化为零，故 \(\int_0^P VIs\,dt=R\int_0^PI^2dt\ge0\)。当 \(P\ll\tau\)，\(h\simeq1-P/(2\tau)\)，切换电流幅值 \(|I_1|\simeq VP/(4L)\)，相对于 \(V/R\) 很小。
